%% kapa: exam=1 round=30 year=2019 period=2 form=A subject=accounting subject_ko=회계학 seq=21 number=61
%% source: Q-Net 감정평가사 기출문제 (cst00309) artlSeq=5205205 file=extracted/1008/5205205/2019년 제30회 감정평가사 1차시험 문제지/2019년 제30회 감정평가사 1차시험 문제지/2019년도 제30회 감정평가사 1차 2교시 A형.hwp
%% conversion: kordoc 4.24.51 -> tools/gen1.py (auto); review=proofread-sonnet; flags=none
\begin{question}{61}
(주)감평은 20x1년 1월 1일에 액면금액 ￦500,000의 전환사채를 다음과 같은 조건으로 액면발행하였다.
\begin{framedtext}
○ 표시이자율: 연 6 \%(매년 말 지급) \par
○ 전환사채 발행당시 일반사채의 시장이자율: 연 10 \% \par
○ 만기일: 20x3년 12월 31일
\end{framedtext}
전환사채의 만기 상환조건이 액면상환조건인 경우의 전환권대가(A)와 할증상환조건(보장수익률 8 \%, 상환할증금 ￦32,464)인 경우의 전환권대가(B)는? (단, 계산금액은 소수점 첫째자리에서 반올림하고, 단수차이로 인한 오차가 있으면 가장 근사치를 선택한다.) \par
{\small\centering
\begin{tabular}{|c|c|c|c|c|}
\hline
\multirow{2}{*}{\makecell{기간}} & \multicolumn{2}{|c|}{\makecell{단일금액 ￦1의 현재가치}} & \multicolumn{2}{|c|}{\makecell{정상연금 ￦1의 현재가치}} \\
\cline{2-5}
 & \makecell{8 \%} & \makecell{10 \%} & \makecell{8 \%} & \makecell{10 \%} \\
\hline
\makecell{3} & \makecell{0.7938} & \makecell{0.7513} & \makecell{2.5771} & \makecell{2.4869} \\
\hline
\end{tabular}\par}

\begin{choices}
  \item A: ￦24,878\hspace{2em}B: ￦488
  \item A: ￦25,787\hspace{2em}B: ￦17
  \item A: ￦25,787\hspace{2em}B: ￦25,353
  \item A: ￦49,743\hspace{2em}B: ￦25,353
  \item A: ￦49,743\hspace{2em}B: ￦17
\end{choices}
\end{question}
\answer{4}
